\documentclass[12pt]{article}
\usepackage[utf8]{inputenc}
\usepackage[slovene]{babel}
\usepackage{amsmath}
\usepackage{graphicx}
\usepackage{float}

\title{3. domača naloga pri predmetu Numerična matematika}
\author{Anja Abramovič}
\date{}

\begin{document}
\maketitle

\section{Opis naloge}
Naloga obravnava reševanje diferencialne enačbe matematičnega nihala:
\begin{equation}
    \frac{g}{l}\sin(\theta(t)) + \ddot{\theta}(t) = 0, \quad \theta(0) = \theta_0, \quad \dot{\theta}(0) = \dot{\theta}_0,
\end{equation}
kjer je $\theta(t)$ kotni odmik, $g$ težni pospešek in $l$ dolžina nihala. Enačbo rešimo numerično z metodo DOPRI5 in primerjamo z analitično rešitvijo harmoničnega nihala. Dodatno izračunamo odvisnost nihajnega časa od energije nihala.


\section{Metoda}
\subsection*{Prevod na sistem prvega reda}
Enačbo drugega reda prevedemo na sistem prvega reda z uvedbo kotne hitrosti 
$\omega = \dot{\theta}$:
\begin{equation}
    \begin{bmatrix} \dot{\theta} \\ \dot{\omega} \end{bmatrix} = 
    \begin{bmatrix} \omega \\ -\frac{g}{l}\sin(\theta) \end{bmatrix}.
\end{equation}

\subsection*{Metoda DOPRI5}
Sistem rešimo z metodo Dormand-Prince reda 5 (DOPRI5), ki je Runge-Kutta metoda z adaptivnim korakom. V vsakem koraku izračunamo 7 vmesnih vrednosti 
$k_1, \ldots, k_7$ in nov približek:
\begin{equation}
    y_{k+1} = y_k + h(b_1 k_1 + b_3 k_3 + b_4 k_4 + b_5 k_5 + b_6 k_6),
\end{equation}
kjer so $b_i$ uteži metode. Napako ocenimo kot razliko med rešitvama reda 4 in 5. Če je napaka prevelika, korak zavržemo in zmanjšamo $h$, sicer pa $h$ prilagodimo:
\begin{equation}
    h_{\text{nov}} = h \cdot \min\left(5, \max\left(0.2, \frac{0.9}{\|e\|^{0.2}}\right)\right).
\end{equation}

\subsection*{Harmonično nihalo}
Za primerjavo uporabimo analitično rešitev harmoničnega nihala 
$\ddot{\theta} + \frac{g}{l}\theta = 0$:
\begin{equation}
    \theta(t) = \theta_0\cos(\omega t) + \frac{\dot{\theta}_0}{\omega}\sin(\omega t),
\end{equation}
kjer je $\omega = \sqrt{g/l}$ kotna frekvenca. Ta aproksimacija je natančna le za majhne odmike.

\subsection*{Nihajni čas}
Nihajni čas matematičnega nihala določimo numerično. Krivuljo $\theta(t)$ rešimo do $2T_\text{harm}$ in poiščemo prvi prehod skozi $\theta = 0$ z bisekcijo. Nihajni čas je štirikratnik tega časa.



\section{Rezultati}
Primerjali smo matematično in harmonično nihalo za dva različna začetna odmika.

Za majhen odmik $\theta_0 = \pi/8$ rad sta rešitvi zelo podobni (slika~\ref{fig:pi8}), saj je pri majhnih kotih $\sin(\theta) \approx \theta$. Za večji odmik $\theta_0 = \pi/2$ rad se rešitvi opazno razlikujeta (slika~\ref{fig:pi2}), matematično nihalo ima daljši nihajni čas ker $\sin(\theta) < \theta$.

\begin{figure}[H]
    \centering
    \includegraphics[width=0.7\textwidth]{primerjava_pi8.png}
    \caption{Primerjava nihala pri $\theta_0 = \pi/8$ rad.}
    \label{fig:pi8}
\end{figure}

\begin{figure}[H]
    \centering
    \includegraphics[width=0.7\textwidth]{primerjava_pi2.png}
    \caption{Primerjava nihala pri $\theta_0 = \pi/2$ rad.}
    \label{fig:pi2}
\end{figure}

Odvisnost nihajnega časa od energije je prikazana na sliki~\ref{fig:nihajni_cas}. Nihajni čas harmoničnega nihala je konstanten ($T \approx 2.006$ s), medtem ko nihajni čas matematičnega nihala narašča z energijo. Ko se odmik približuje $\pi$, gre nihajni čas proti neskončnosti.

\begin{figure}[H]
    \centering
    \includegraphics[width=0.7\textwidth]{nihajni_cas.png}
    \caption{Odvisnost nihajnega časa matematičnega nihala od energije.}
    \label{fig:nihajni_cas}
\end{figure}





\end{document}